\section{Una corrida frente a repeticiones}\label{sec:variabilidad}

Las dos baterías ejecutaron cada configuración una sola vez; la campaña repitió algunas de ellas de tres a cinco veces. Donde coinciden la prueba, el tamaño, la red y el número de procesos, se puede comparar la magnitud de las diferencias entre campañas. Como también cambió la compilación, esas diferencias no estiman por sí solas el error de una corrida única. La tabla usa la variante 1D global de la campaña y, para el trapecio, su modo simétrico.

\begin{table}[htbp]
\caption{Tiempo total (s) de las corridas únicas de cada batería frente a la mediana y el rango de la campaña.}\label{tab:variabilidad}
\begin{center}\small
\begin{tabular}{llrrrrrr}
\toprule
 & & & \multicolumn{2}{c}{Corrida única} & \multicolumn{3}{c}{Campaña} \\
\cmidrule(lr){4-5}\cmidrule(lr){6-8}
Prueba & Red & $P$ & Bat. 1 & Bat. 2 & Mediana & Rango & Reps. \\
\midrule
Matrices $N=3072$ & Ethernet & 24 & 0.512 & 0.547 & 0.649 & 0.633--0.719 & 5 \\
Matrices $N=3072$ & Ethernet & 96 & 6.175 & 6.148 & 6.260 & 6.233--6.266 & 5 \\
Matrices $N=3072$ & Wi-Fi & 24 & 0.553 & 0.553 & 0.666 & 0.581--0.722 & 5 \\
Matrices $N=3072$ & Wi-Fi & 96 & 101.042 & 99.945 & 101.082 & 100.293--102.991 & 5 \\
Matrices $N=6144$ & Ethernet & 24 & -- & 3.839 & 5.279 & 4.998--7.343 & 3 \\
Matrices $N=6144$ & Ethernet & 96 & -- & 25.037 & 25.194 & 25.065--25.322 & 3 \\
Trapecio $n=10^{11}$ & Ethernet & 24 & 3.870 & 3.842 & 5.183 & 4.753--5.379 & 3 \\
Trapecio $n=10^{11}$ & Ethernet & 96 & 0.964 & 0.975 & 1.012 & 1.004--1.024 & 3 \\
\bottomrule
\end{tabular}
\end{center}

\end{table}

\subsection{Qué muestra la comparación}

\begin{itemize}
\item \textbf{Con 96 procesos, las corridas únicas de matrices quedaron próximas a las medianas posteriores.} Con $N=3072$ y $N=6144$, las baterías quedan a menos de un 2\,\% de la mediana repetida en las redes medidas. El tiempo está dominado por la comunicación, que dio resultados próximos en los tres momentos.
\item \textbf{Con 24 procesos, la campaña fue sistemáticamente más lenta:} un 19--27\,\% con $N=3072$ y un 37\,\% con $N=6144$. Ninguna de las corridas repetidas de $N=3072$ por Ethernet bajó hasta los tiempos de las baterías.
\item \textbf{La mayor diferencia observada está en el cálculo.} La tabla~\ref{tab:variabilidad-fases} separa las fases de las configuraciones con 24 procesos.
\end{itemize}

\begin{table}[htbp]
\caption{Máximos de cálculo y de distribución con 24 procesos: corridas únicas frente a la mediana de la campaña.}\label{tab:variabilidad-fases}
\begin{center}\small
\begin{tabular}{llrrrrrr}
\toprule
 & & \multicolumn{3}{c}{$T_{Calc}$ máximo (s)} & \multicolumn{3}{c}{$T_{Dist}$ máximo (s)} \\
\cmidrule(lr){3-5}\cmidrule(lr){6-8}
Matriz & Red & Bat. 1 & Bat. 2 & Campaña & Bat. 1 & Bat. 2 & Campaña \\
\midrule
$N=3072$ & Ethernet & 0.304 & 0.326 & 0.455 & 0.198 & 0.205 & 0.220 \\
$N=3072$ & Wi-Fi & 0.345 & 0.337 & 0.450 & 0.216 & 0.228 & 0.214 \\
$N=6144$ & Ethernet & -- & 2.891 & 4.447 & -- & 0.875 & 0.959 \\
\bottomrule
\end{tabular}
\end{center}

\end{table}

$T_{Dist}$ cambia poco (hasta un 11\,\%), mientras $T_{Calc}$ crece entre 1.3 y 1.5 veces en la campaña (0.304 a 0.455 s con $N=3072$ por Ethernet; 2.891 a 4.447 s con $N=6144$). La diferencia es compatible con el cambio de compilación: las baterías registran \texttt{-O3 -march=native -funroll-loops}, mientras la campaña registra \texttt{-O3 -std=c11 -Wall -Wextra}. No se inspeccionó el ensamblado ni se conservaron informes de vectorización que permitan atribuirla a instrucciones concretas. También cambió el momento de ejecución. Una prueba causal requiere recompilar el mismo código, alternar binarios bajo condiciones compartidas y repetir.

El trapecio muestra un patrón similar con 24 procesos (3.84--3.87 s en las baterías frente a 5.18 s de mediana), pero ahí la comparación mezcla además dos programas: las baterías usaron \archivo{trapecio.c} y la campaña, \archivo{trapecio_asimetrico} en modo simétrico, con la corrección del extremo.

\subsection{Consecuencias para la lectura del informe}

\begin{enumerate}
\item Las conclusiones cualitativas se mantienen: en las dos baterías, un nodo con 24 procesos es el mejor punto de los valores de $P$ ensayados para matrices y la difusión entre nodos domina el tiempo con 96 procesos.
\item Los cocientes $T_{96}/T_{24}$ de la campaña (4.77 y 4.57) no deben mezclarse con los de las baterías (6.52 para $N=6144$). Si mejorara únicamente el cálculo manteniendo iguales los demás costes, la configuración local, más sensible a ese cálculo, podría ganar más; los datos actuales no aíslan esa intervención.
\item Una diferencia pequeña entre dos corridas únicas con 24 procesos, por ejemplo entre redes, no es interpretable: la variación entre momentos supera esas diferencias.
\end{enumerate}
